\subsection*{一步一步建立直观}

\paragraph{1. 先读位置、速度，再读速度的变化}

一次导数描述输出对输入的小变化有多敏感，二次导数则描述这种敏感度如何改变。把输入看成时间时，位置、速度、加速度依次出现。右图用指数函数的多项式近似演示逐阶匹配：一次项调整倾斜，二次项开始调整弯曲。



\paragraph{2. 同一条切线不能规定整条曲线}

把观察位置移得更远，低阶多项式与原函数可能迅速分离。两条曲线在一点共享函数值与一阶导数，只保证它们在一阶尺度上接近；二阶导数及更高阶信息可以不同。因此切线相同不是函数相同，更不是远处预测可靠。



\paragraph{3. 高阶信息增加局部约束}

增加阶数，检查读数中的实际误差。这里指数函数在零点的各阶导数都等于一，所以系数写成阶乘分母。一般函数需要分别计算各阶导数，并估计余项。即使已匹配许多阶，仍应同时注明近似中心、使用区间和允许误差。
